% Proposed replacement for the Response to Referee 1, Comment 1.
% Author review, 24 October 2026; NOT automatically approved for submission.
% Based on the supplied r1-1.pdf and ref1_reply.pdf.
% The numerical section explicitly uses the independently completed current
% LO + quark + leading mixed-condensate OPE. It is NOT the original-paper rerun.
% Before submission, discuss this scope and consistently implement/recompute
% the outstanding Appendix C correction and its dependent results.
% No manuscript file or baseline numerical output has been changed.

\noindent\textbf{Response:}

We thank the referee for this helpful suggestion. We have expanded
Sec.~II, following the two-point sum rules for the decay constants, to
summarize the mixing-angle determination in our earlier work. The added
discussion, Eqs.~(39)--(43) of the revised manuscript, explains the
current rotation, the diagonalization condition, and the role of the
physical diagonal correlators. The published mixing angles are quoted
in Eq.~(46).

Within the adopted two-state description, the angle is determined by
requiring the Borel-transformed and continuum-subtracted spin-one
correlation matrix to be diagonal in the rotated current basis.
For the current convention used in the manuscript,
\begin{equation}
R(\theta)=
\begin{pmatrix}
\sin\theta & \cos\theta\\
\cos\theta & -\sin\theta
\end{pmatrix},
\qquad
\left[R(\theta)\widehat{\boldsymbol\Pi}R^T(\theta)\right]_{12}=0.
\end{equation}
With a common Borel parameter and continuum-subtraction prescription,
this gives
\begin{equation}
\tan(2\theta)=-\frac{2\widehat\Pi_{AB}(M^2,s_0)}
{\widehat\Pi_{AA}(M^2,s_0)-\widehat\Pi_{BB}(M^2,s_0)}.
\end{equation}
The two physical diagonal correlators need not be equal: their
difference cancels in the ratio wherever it is defined. Different
physical masses and pole residues are therefore compatible with this
relation within the adopted description. The earlier results,
\begin{equation}
\theta_{D_s}=(26.6\pm0.6)^\circ,
\qquad
\theta_{B_s}=(38.5\pm0.1)^\circ,
\end{equation}
are retained, with their quoted uncertainties, as external inputs.

We have also added a footnote clarifying a typographical omission in
the earlier article: the tensor-current normalization factor
$1/(m_Q+m_s)$, with $Q=c,b$, was included in the calculations but
inadvertently omitted from the printed current definition. The two
studies thus use the same current normalization.

The light-quark condensate at $1\,\mathrm{GeV}$ changes from
$-(1.65\pm0.15)\times10^{-2}\,\mathrm{GeV}^3$ in the earlier study
to $-(0.240\pm0.010)^3\,\mathrm{GeV}^3$ in the present analysis.
The central value of $\langle\bar ss\rangle/\langle\bar qq\rangle$
remains $0.80$, while its uncertainty changes from $0.20$ to $0.10$.

To assess the sensitivity to the central condensate value, we performed
an independent paired check of the mixing-angle ratio using the
leading-order perturbative contribution, the quark-condensate terms,
and the complete leading dimension-five mixed-condensate contribution.
Gluon-condensate and radiative corrections are omitted in this
auxiliary check. With this OPE prescription and all other inputs fixed,
we changed
\begin{equation}
\langle\bar qq\rangle=-0.0165\,\mathrm{GeV}^3
\quad\longrightarrow\quad
-0.013824\,\mathrm{GeV}^3
\qquad (\mu=1\,\mathrm{GeV}),
\end{equation}
changing the strange and mixed condensates consistently through
\begin{equation}
\langle\bar ss\rangle=0.8\langle\bar qq\rangle,
\qquad
\langle\bar s g_s\sigma\!\cdot\!G s\rangle
=m_0^2\langle\bar ss\rangle,
\qquad m_0^2=0.8\,\mathrm{GeV}^2.
\end{equation}
At $(M^2,s_0)=(4.5,10)\,\mathrm{GeV}^2$ for $D_s$ and
$(11,43)\,\mathrm{GeV}^2$ for $B_s$, the shifts
$\Delta\theta=\theta_{\mathrm{new}}-\theta_{\mathrm{old}}$ are
\begin{equation}
\Delta\theta_{D_s}=+0.175^\circ,
\qquad
\Delta\theta_{B_s}=+0.045^\circ.
\end{equation}
A $21\times21$ grid gives
\begin{equation}
\Delta\theta_{D_s}=0.116^\circ\text{--}0.331^\circ,
\qquad
\Delta\theta_{B_s}=0.036^\circ\text{--}0.060^\circ,
\end{equation}
over $M^2=3$--$6\,\mathrm{GeV}^2$, $s_0=9$--$11\,\mathrm{GeV}^2$
for $D_s$, and $M^2=9$--$13\,\mathrm{GeV}^2$,
$s_0=42$--$44\,\mathrm{GeV}^2$ for $B_s$.
Within this specified truncation, the shifts remain below the quoted
angle uncertainties throughout the sampled regions. This is a
central-input sensitivity test, not a reproduction of the full earlier
extraction or a new propagation of its uncertainties; in particular,
the changed uncertainty in the strange-to-light condensate ratio is
not propagated here. The published mixing angles remain the external
inputs of the radiative-decay analysis.
